% Generated from tex/chapters/ch04/08_構築技術.tex for the SWEBOK structure tree
\subsubsection{エラー処理、例外処理、フォールトトレランス}

エラー処理の方法は、正しさ、頑健性、信頼性に影響する。エラー処理には、警告ログを出す、エラーコードを返す、中立値を返す、次の有効データを使う、処理を停止するなど、複数の方法がある。どれを使うかは、利用者への影響、データの重要性、安全性、復旧可能性によって変わる。

例外処理は、エラーや例外的な出来事を検出し、処理する仕組みである。基本的には、例外を送出し、対応する処理ブロックで捕捉する。例外メッセージには、原因理解に必要な情報を含めるべきである。空の捕捉ブロックで例外を握りつぶすと、問題の発見が遅れる。

フォールトトレランスとは、エラーを検出し、回復するか、影響を閉じ込めることで、ソフトウェアの信頼性を高める技法群である。再試行、補助コード、投票アルゴリズム、値の代替などがある。
